% section: 漸化式の分類

\section{漸化式の分類}\label{節：漸化式の分類}
  \BookFigClassification

  ここまでに使った方法を、一枚の地図として見直す。
  まず\bookterm{recurrence}{漸化式}は解ける\bookterm{recurrence}{漸化式}と解けない\bookterm{recurrence}{漸化式}に分類される.
  そして解ける\bookterm{recurrence}{漸化式}は「決まった手順で解ける\bookterm{recurrence}{漸化式}」と「特定の性質から解ける\bookterm{recurrence}{漸化式}」に分類される.
  「決まった手順で解ける\bookterm{recurrence}{漸化式}」はさらに分類され,各分類ごとに様々な解法がある.
  ここで言う「解ける」というのは,本書内での定義であり,一般的なものではない.
  ここで解けないとして扱われるような\bookterm{recurrence}{漸化式}は\bookterm{general}{一般項}を表す手段が存在しないわけではない.
  あくまで本書で扱う手法の内では表現できないという意味である.
  「解ける」の定義について述べておく.
  定義は,高校数学程度の範囲内で既知の関数に帰着するか,\,$n$に依存しない有限個の\bookterm{elementary}{初等関数}の演算によって\bookterm{general}{一般項}を表せるような解法が存在するものを意味している.
  このうち「決まった手順で解ける\bookterm{recurrence}{漸化式}」は基本4パターン・応用7パターン・発展パターンの各節で,「特定の性質から解ける\bookterm{recurrence}{漸化式}」は特殊性質の節で扱う.

  基本4パターンでは、差と倍率から項を求めた。
  応用パターンでは、追う量を変えることで、その基本形を取り出した。
  発展パターンと特殊性質型では、追う対象と調べる性質がさらに広がった。

  この分類は,解法を探すための案内である.一つの式に複数の解法が使えることもある.
  本書の方法で\bookterm{closed}{閉じた形}が見つからなくても,\bookterm{sequence}{数列}が定義されないわけではなく,\bookterm{limit}{極限}や整数の性質を調べる価値も失われない.
  例えば\bookproblemref{practice-entrance}{6}{発展問6}では,\bookterm{general}{一般項}を求めずに減り方を調べる.
